\section{Introduction}

Differentiable logic gate networks (DLGNs) train Boolean circuits by optimizing continuous gate relaxations and discretizing the learned functions for inference \cite{petersen2022}. Their restricted operators make the connection graph particularly consequential: a two-input gate can combine only the two signals assigned to it. Gate training cannot introduce an absent predecessor, making connectivity a design constraint at a given depth and gate budget.

Existing approaches make this decision either before or during training. Fixed random connections retain a simple graph representation, while learned routing optimizes the predecessors together with the gate functions \cite{mommen2025,fojcik2026}. Learned routing can improve accuracy, but processing candidate inputs increases the training state even when the final circuit retains only two connections per gate. In our local comparison under two LILogic architectures, Top-32 routing requires 16.0–17.1 times the peak GPU memory of fixed connectivity. This cost motivates a complementary question: how much accuracy can a deliberately designed fixed graph recover before routing itself must be learned?

The relevant structural property extends beyond balanced predecessor use. The native dense random baseline already balances fan-out, yet graphs with the same predecessor degrees need not expose the same combinations of image features. Our dense ablation makes this distinction concrete: randomizing first-layer connections while preserving per-slot predecessor degrees reduces mean validation accuracy by 4.49 percentage points with structured deeper layers. Thus, the identity of the paired signals can matter even when local degree statistics remain unchanged.

We introduce LogicConnect to exploit this design opportunity through semantic input pairing and regular multiscale matchings. Semantic pairs preserve the relationship between encoded bits and their underlying image sources. Deeper matchings organize signal combinations across index scales while controlling predecessor reuse. An ancestry-overlap criterion chooses among equally balanced stages without learning routing parameters. Applying these principles to convolutional DLGNs requires a channel-level construction that preserves shared logic trees and spatial sampling, rather than treating every spatial gate application as an independent dense gate.

Our contributions are threefold:

\begin{itemize}
\item We introduce a unified fixed-connectivity construction for dense gates, convolutional channel groups, and their classifiers. The construction runs before training and preserves the gate budget, gate parameterization, and inference operators.
\item We show that LogicConnect improves dense CIFAR-10 test accuracy by 3.18–4.59 percentage points across three sizes and three paired seeds. We characterize its accuracy/resource trade-off against learned routing and its application to convolutional models.
\item We isolate structural effects with degree-preserving randomizations and simpler stage-selection controls. These experiments identify the dense first-layer contribution and bound the claims about deeper structure, ancestry selection, and transfer.

\end{itemize}
Convolutional S/M models improve by 3.26/2.08 points in existing single-seed tests, while the three-seed S validation interval includes zero. Learned routing remains more accurate in the matched comparison, establishing where additional training memory still buys accuracy.
